\chapter{Complexity — V Word}
%%% generated by the blueprint pipeline from TCSlib.Complexity.SpaceComplexity.Machines.VWord
\section{Overview}

Declarations of the Lean module \texttt{TCSlib.Complexity.SpaceComplexity.Machines.VWord}.

\section{Declarations}

\begin{lemma}[Indexing after dropping a list prefix]
\lean{Complexity.getD_drop}
\label{Complexity.getD_drop}
\leanok
\uses{List.getD_drop}
\difficulty{1}
For every Boolean list $y$ and natural numbers $a,i$, the element at index $i$ of the
list obtained by dropping the first $a$ entries of $y$, using false as the default when
the index is out of range, equals the element of $y$ at index $a+i$ with the same
default.
\end{lemma}

\begin{lemma}[List entry after dropping a prefix]
\lean{Complexity.getD_of_drop}
\label{Complexity.getD_of_drop}
\leanok
\uses{Complexity.getD_drop}
\difficulty{3}
Let $x$ and $r$ be Boolean lists, let $p$ be a natural number, and let $b$ be a Boolean.
If dropping the first $p$ entries of $x$ yields the list whose first entry is $b$ and
whose remaining entries form $r$, then the entry of $x$ at index $p$, using false as the
default for an out-of-range index, is $b$.
\end{lemma}

\begin{lemma}[Successive list drop]
\lean{Complexity.drop_succ_of_drop}
\label{Complexity.drop_succ_of_drop}
\leanok
\difficulty{3}
Let $x$ and $r$ be Boolean lists, let $p$ be a natural number, and let $b$ be a Boolean.
If dropping the first $p$ entries of $x$ gives the list whose first entry is $b$ and
whose remaining list is $r$, then dropping the first $p+1$ entries of $x$ gives $r$.
\end{lemma}

\begin{lemma}[List Drop Bound]
\lean{Complexity.lt_of_drop}
\label{Complexity.lt_of_drop}
\leanok
\difficulty{4}
For lists $x,r$ of Boolean values, a natural number $p$, and a Boolean value $b$, if
deleting the first $p$ entries of $x$ gives a list whose first entry is $b$ and whose
remaining tail is $r$, then $p<|x|$.
\end{lemma}

\begin{lemma}[Successive list drops]
\lean{Complexity.drop_add_of_drop}
\label{Complexity.drop_add_of_drop}
\leanok
\difficulty{3}
Let $x,u,r$ be Boolean lists and let $p$ be a natural number. If dropping the first $p$
entries of $x$ leaves the concatenation of $u$ and $r$, then dropping the first $p+|u|$
entries of $x$ leaves $r$.
\end{lemma}

\begin{lemma}[Length bound after dropping a prefix]
\lean{Complexity.le_of_drop}
\label{Complexity.le_of_drop}
\leanok
\difficulty{4}
Let $x,u,r$ be finite binary strings and let $p$ be a natural number. If deleting the
first $p$ entries of $x$ leaves the concatenation of $u$ and $r$, and $u$ is nonempty,
then $p+|u|\le |x|$.
\end{lemma}

\begin{theorem}[Whole input with two registers]
\lean{Complexity.LogProg.vword_whole₂}
\leanok
\uses{Complexity.LogProg.Mode.seg0, Complexity.LogProg.argSegs, Complexity.LogProg.callSegs, Complexity.LogProg.render, Complexity.LogProg.vword, Turing.pairEncode, Turing.pairEncode_eq_dbl}
\difficulty{1}
For any bit string $u$, decider index $j$, register indices $r,s$, states $l_1,l_0$, and
assignment $V$ of a bit string to each register, the virtual word formed using the
whole-input mode on $u$ and the two registers $r,s$ is the encoded pair $(u,	ext{the
encoded pair }(V(r),V(s)))$.
\end{theorem}

\begin{theorem}[Whole input with three registers]
\lean{Complexity.LogProg.vword_whole₃}
\leanok
\uses{Complexity.LogProg.Mode.seg0, Complexity.LogProg.argSegs, Complexity.LogProg.callSegs, Complexity.LogProg.render, Complexity.LogProg.vword, Turing.pairEncode, Turing.pairEncode_eq_dbl}
\difficulty{1}
For any bit string $u$, decider index $j$, register indices $r,s,t$, states $l_1,l_0$,
and assignment $V$ of a bit string to each register, the virtual word formed using the
whole-input mode on $u$ and the three registers $r,s,t$ is the encoded pair $(u,	ext{the
encoded pair }(V(r),	ext{the encoded pair }(V(s),V(t))))$.
\end{theorem}
